\documentclass[version=last, paper=letter, fromalign=right, fromphone,
  fromemail, foldmarks=false, parskip=half]{scrlttr2}
\LoadLetterOption{UScommercial9}
\usepackage[T1]{fontenc}
\usepackage[french]{babel}

\setkomavar{fromname}{{{AUTHOR}}}
\setkomavar{fromaddress}{123, rue Principale\\Ville (Province)\\A1B 2C3}
\setkomavar{fromphone}{514 555-0123}
\setkomavar*{fromphone}{Téléphone~:}
\setkomavar{fromemail}{prenom.nom@exemple.org}
\setkomavar*{fromemail}{Courriel~:}
\setkomavar{place}{Ville}
\setkomavar{date}{\today}
\setkomavar{subject}{Objet de la lettre}
\setkomavar{signature}{{{AUTHOR}}}

\begin{document}

\begin{letter}{%
  Nom du destinataire, fonction\\
  Organisme\\
  456, avenue des Érables\\
  Ville (Province)\\
  D4E 5F6%
}

\opening{Madame, Monsieur,}

Le premier paragraphe annonce l'objet de la lettre. Les suivants en
développent le contenu, un sujet par paragraphe.

Le dernier paragraphe indique la suite attendue: une réponse, une rencontre
ou un document à fournir.

\closing{Veuillez agréer, Madame, Monsieur, mes salutations distinguées.}

\encl{Curriculum vitæ}

\end{letter}

\end{document}
